\chapter*{Conclusion générale}
\label{chap:conclusion}
\addstarredchapter{Conclusion générale}

\initial{C}e stage posait une question précise : comment estimer, sans fuite et avec une validation externe, un risque de conversion MCI $\rightarrow$ démence à 2 et 5 ans, et quel paquet de données minimal reste utile. La littérature montre que l'AUC interne sur ADNI n'est plus une contribution. Le travail livré est un protocole aligné sur le canevas du laboratoire Fongang, une cartographie des cohortes NACC, ADNI et Framingham, et une revue qui fixe le langage du labo (CARD, Kolachalama, MOIRA, PROBAST).

La démarche -- phénotype locké, split participant, profils incomplets conservés, modèles emboîtés, calibration et \ac{DCA} -- est la condition pour que les chiffres du mois 2, lorsqu'ils seront produits sous DUA, soient interprétables. Les limites (durée, représentativité, données africaines hors fenêtre) sont énoncées plutôt que dissimulées. L'intérêt de ce stage, au-delà d'un modèle, est d'avoir rendu le projet \emph{exécutable} dans les règles d'un laboratoire de neurosciences computationnelles.
